% Syntra LaTeX Template
% For Arabic support, use XeLaTeX with:
%   \usepackage{fontspec}
%   \usepackage{polyglossia}
%   \setdefaultlanguage{english}
%   \setotherlanguage{arabic}
%   \newfontfamily\arabicfont[Script=Arabic]{Amiri}

\documentclass[12pt,a4paper]{article}

% Font and encoding
\usepackage[utf8]{inputenc}
\usepackage[T1]{fontenc}

% Math packages
\usepackage{amsmath,amssymb}

% Table support
\usepackage{longtable,booktabs}

% Graphics and links
\usepackage{graphicx}
\usepackage{hyperref}
\usepackage{geometry}

% Page geometry
\geometry{
    left=2.5cm,
    right=2.5cm,
    top=3cm,
    bottom=3cm
}

% Fix pandoc tightlist issue
\providecommand{\tightlist}{\setlength{\itemsep}{0pt}\setlength{\parskip}{0pt}}

% Hyperref settings
\hypersetup{
    colorlinks=true,
    linkcolor=blue,
    filecolor=magenta,
    urlcolor=cyan,
    pdftitle={Syntra Research Documentation},
    pdfauthor={Syntra-Env},
    pdfpagemode=FullScreen
}

\begin{document}

\section{Scholar Mathematical
Framework}\label{scholar-mathematical-framework}

Scholar uses three mathematical frameworks to analyze the Quran
computationally: a \emph{distributional root space} for measuring how
Arabic roots behave across the text, an \emph{su(2) gauge field} for
detecting structural tension and meaning shifts, and a \emph{UOR ring
substrate} that provides content-addressed identity for all textual
units. All math lives in \texttt{src/math/} as pure functions with no
side effects.

\subsection{Contents}\label{contents}

\begin{enumerate}
\tightlist
\item
  \hyperref[root-space]{Root Space (Distributional Vectors)}
\item
  \hyperref[verse-dynamics]{Verse Dynamics}
\item
  \hyperref[gauge]{su(2) Gauge Field}
\item
  \hyperref[uor]{UOR Ring Substrate}
\item
  \hyperref[integration]{Integration}
\end{enumerate}

\subsection{1. Root Space (Distributional Vectors)}\label{root-space}

Every Arabic root receives a \textbf{127-dimensional distributional
fingerprint} built from its grammatical usage across the entire Quran.

\subsubsection{1.1 Feature Vocabulary}\label{feature-vocabulary}

Category

Dims

What it captures

POS tags

44

Noun, verb, particle, \ldots{}

Verb forms

12

Form I--X + derived patterns

Aspect

3

Perfect, imperfect, imperative

Person

3

1st, 2nd, 3rd

Number

3

Singular, dual, plural

Gender

2

Masculine, feminine

Case

3

Nominative, accusative, genitive

\textbf{Total}

\textbf{127}

Table 1. Feature dimensions of the root profile space.

The profile is a normalized probability distribution: for each root,
count how often it appears in each feature slot, then normalize so the
entries sum to 1.

\subsubsection{1.2 Distributional Weight}\label{distributional-weight}

Measures the structural importance of a root:

\$\$w =
0.6\textbackslash;\textbackslash frac\{\textbar\textbackslash text\{surahs\}\textbar\}\{114\}
+
0.4\textbackslash;\textbackslash frac\{\textbar\textbackslash text\{morpheme
types\}\textbar\}\{\textbar\textbackslash text\{instances\}\textbar\}
\textbackslash tag\{1\}\$\$

where the first term is \emph{surah spread} and the second is \emph{form
diversity}.

\subsubsection{1.3 Instance Anomaly}\label{instance-anomaly}

Scores how unusual a specific usage of a root is (0 = typical, 1 =
never-before-seen form):

\$\$a\_i = 1 - p{[}f\_i{]} \textbackslash tag\{2\}\$\$

where \$p\$ is the root\textquotesingle s profile distribution and
\$f\_i\$ is the feature index of the instance. A root that usually
appears as a noun but occurs as a verb scores high anomaly on that
instance.

\subsubsection{1.4 Concordance Distance}\label{concordance-distance}

The distance between two roots \$A\$ and \$B\$ is the Jensen--Shannon
divergence of their profiles:

\$\$D\_\{\textbackslash text\{JS\}\}(A \textbackslash\textbar{} B) =
\textbackslash frac\{1\}\{2\}\textbackslash,D\_\{\textbackslash text\{KL\}\}\textbackslash!\textbackslash left(p\_A
\textbackslash,\textbackslash Big\textbackslash\textbar\textbackslash,
\textbackslash frac\{p\_A + p\_B\}\{2\}\textbackslash right) +
\textbackslash frac\{1\}\{2\}\textbackslash,D\_\{\textbackslash text\{KL\}\}\textbackslash!\textbackslash left(p\_B
\textbackslash,\textbackslash Big\textbackslash\textbar\textbackslash,
\textbackslash frac\{p\_A + p\_B\}\{2\}\textbackslash right)
\textbackslash tag\{3\}\$\$

This measures how differently two roots behave grammatically.

\subsection{2. Verse Dynamics}\label{verse-dynamics}

Verse analysis computes four signals per word, then aggregates them into
an emphasis profile for the verse.

\subsubsection{2.1 Per-Word Signals}\label{per-word-signals}

Signal

Symbol

Computation

Weight

\$w\_i\$

Distributional weight of the word\textquotesingle s root (Eq.~1)

Anomaly

\$a\_i\$

Instance anomaly score (Eq.~2)

Drift

\$\textbackslash delta\_i\$

Profile distance between consecutive root vectors

Coherence

\$c\_i\$

Profile cosine similarity with the next root-bearing word

Table 2. Four signals computed for each word in a verse.

\subsubsection{2.2 Emphasis Profile}\label{emphasis-profile}

The four per-word signals aggregate into a verse-level emphasis
breakdown:

\begin{itemize}
\tightlist
\item
  \textbf{Root emphasis} --- dominated by high-weight words
  (structurally important roots)
\item
  \textbf{Morpheme emphasis} --- dominated by high-anomaly words
  (unusual grammatical forms)
\item
  \textbf{Position emphasis} --- dominated by high-drift words (meaning
  shifts at specific positions)
\item
  \textbf{Boundary signal} --- inverse of minimum coherence (low
  coherence ⇒ boundary)
\end{itemize}

The dominant emphasis type tells you \emph{what kind of meaning} the
verse carries: structural (root), exceptional (morpheme), positional
(drift), or transitional (boundary).

\subsubsection{2.3 Verse Coherence
(Phase-Lock)}\label{verse-coherence-phase-lock}

Measures whether two verses ``belong together''---a value in
\${[}0,1{]}\$:

\$\$\textbackslash Phi =
\textbackslash frac\{1\}\{3\}\textbackslash bigl(\textbackslash text\{instance\textbackslash\_agreement\}
+ \textbackslash text\{root\textbackslash\_drift\} +
\textbackslash text\{topic\textbackslash\_overlap\}\textbackslash bigr)
\textbackslash tag\{4\}\$\$

\begin{itemize}
\tightlist
\item
  \textbf{Instance agreement}: average cosine similarity of shared
  roots\textquotesingle{} instance profiles
\item
  \textbf{Root drift}: shift in the distributional landscape between
  verses
\item
  \textbf{Topic overlap}: Jaccard overlap of active roots
\end{itemize}

Thresholds: \$\textbackslash Phi \textgreater{} 0.5\$ = continuation,
\$0.25 \textless{} \textbackslash Phi \textless{} 0.5\$ = transition,
\$\textbackslash Phi \textless{} 0.25\$ = hard boundary.

\subsection{3. su(2) Gauge Field}\label{gauge}

The \textbf{abjadic gauge field} maps Quranic morphological features
onto the \$\textbackslash mathfrak\{su\}(2)\$ Lie algebra---the same
algebra that describes spin-\$\textbackslash tfrac\{1\}\{2\}\$ particles
in physics. The mathematical structure (non-commutative gauge
connections, holonomy, curvature) is used literally, not as metaphor.

\subsubsection{3.1 Pauli Basis}\label{pauli-basis}

The \$\textbackslash mathfrak\{su\}(2)\$ generators are the Pauli
matrices:

\$\$\textbackslash sigma\_1 = \textbackslash begin\{pmatrix\} 0 \& 1
\textbackslash\textbackslash{} 1 \& 0
\textbackslash end\{pmatrix\},\textbackslash quad
\textbackslash sigma\_2 = \textbackslash begin\{pmatrix\} 0 \& -i
\textbackslash\textbackslash{} i \& 0
\textbackslash end\{pmatrix\},\textbackslash quad
\textbackslash sigma\_3 = \textbackslash begin\{pmatrix\} 1 \& 0
\textbackslash\textbackslash{} 0 \& -1 \textbackslash end\{pmatrix\}
\textbackslash tag\{5\}\$\$

\subsubsection{3.2 Feature-to-su(2)
Mapping}\label{feature-to-su2-mapping}

Each morpheme instance maps to a point in
\$\textbackslash mathfrak\{su\}(2)\$ via three components:

\$\$H = x\textbackslash,\textbackslash sigma\_1 +
y\textbackslash,\textbackslash sigma\_2 +
z\textbackslash,\textbackslash sigma\_3 \textbackslash tag\{6\}\$\$

Component

Axis

Features mapped

\$x\$

Semantic Identity

Root surprisal + lemma surprisal + UOR coupling

\$y\$

Morphological Action

Verb form, aspect, person, gender, number surprisals

\$z\$

Syntactic Position

POS, case, voice, mood surprisals

Table 3. Mapping of linguistic features to
\$\textbackslash mathfrak\{su\}(2)\$ axes.

\textbf{Surprisal} is the information-theoretic content of a feature:

\$\$s(f) =
\textbackslash log\textbackslash!\textbackslash left(\textbackslash frac\{N\}\{n\_f\}\textbackslash right)
\textbackslash tag\{7\}\$\$

where \$N = 128\{,\}000\$ (corpus scale) and \$n\_f\$ is the
feature\textquotesingle s frequency. Rare features contribute more
signal. The resulting vector \$(x,y,z)\$ is normalized to the unit
sphere before constructing \$H\$.

\subsubsection{3.3 Curvature Tensor}\label{curvature-tensor}

For two connections \$H\_\textbackslash mu\$ and
\$H\_\textbackslash nu\$, the curvature is the Lie bracket:

\$\$R\_\{\textbackslash mu\textbackslash nu\} =
{[}H\_\textbackslash mu,\textbackslash, H\_\textbackslash nu{]} =
H\_\textbackslash mu H\_\textbackslash nu - H\_\textbackslash nu
H\_\textbackslash mu \textbackslash tag\{8\}\$\$

Non-zero curvature means the order of traversal matters---the gauge
field has genuine structure. The scalar field tension is:

\$\$\textbackslash mathcal\{T\} =
\textbackslash operatorname\{Tr\}(R\_\{\textbackslash mu\textbackslash nu\}\^{}\textbackslash dagger\textbackslash,
R\_\{\textbackslash mu\textbackslash nu\}) \textbackslash tag\{9\}\$\$

\subsubsection{3.4 Discrete Holonomy}\label{discrete-holonomy}

For a sequence of connections \$(H\_1, H\_2, \textbackslash ldots,
H\_n)\$---e.g. all instances of a root---compute the path-ordered
product:

\$\$U = H\_n \textbackslash cdot H\_\{n-1\} \textbackslash cdots H\_1
\textbackslash tag\{10\}\$\$

The holonomy angle is extracted via:

\$\$\textbackslash kappa =
2\textbackslash,\textbackslash arccos\textbackslash!\textbackslash left(\textbackslash left\textbar\textbackslash frac\{\textbackslash operatorname\{Tr\}(U)\}\{2\}\textbackslash right\textbar\textbackslash right)
\textbackslash tag\{11\}\$\$

This is the total ``rotation'' accumulated by traversing all instances
of a root.

\subsubsection{3.5 Root Resonance}\label{root-resonance}

The holonomy angle \$\textbackslash kappa\$ applied across all instances
of a root characterizes its semantic behaviour:

\begin{itemize}
\tightlist
\item
  \$\textbackslash kappa \textbackslash to 0\$: meaning is
  \emph{conserved} (consistent usage everywhere)
\item
  \$\textbackslash kappa \textbackslash to \textbackslash pi\$: meaning
  \emph{evolves} (highly context-dependent)
\end{itemize}

Validated examples:

\begin{itemize}
\tightlist
\item
  {كتب} (k-t-b): \$\textbackslash kappa = 2.35\$ --- context-dependent
\item
  {ربب} (r-b-b): \$\textbackslash kappa = 1.06\$ --- moderately
  consistent
\end{itemize}

\subsubsection{3.6 Lyapunov Deviation}\label{lyapunov-deviation}

Measures how far a specific instance deviates from the
root\textquotesingle s average gauge state:

\$\$L = \textbackslash left\textbackslash\textbar H -
H\^{}*\textbackslash right\textbackslash\textbar\_F\^{}2 =
\textbackslash operatorname\{Tr\}\textbackslash!\textbackslash bigl((H -
H\^{}*)\^{}\textbackslash dagger (H - H\^{}*)\textbackslash bigr)
\textbackslash tag\{12\}\$\$

where \$H\^{}* =
\textbackslash frac\{1\}\{n\}\textbackslash sum\_\{i=1\}\^{}n H\_i\$ is
the mean connection matrix. High-deviation instances are potential
falsifiers---places where a proposed meaning for the root might not
hold.

\subsection{4. UOR Ring Substrate}\label{uor}

The Universal Object Reference (UOR) provides content-addressed identity
for every textual unit. All arithmetic happens in the ring
\$\textbackslash mathbb\{Z\}/2\^{}\{256\}\textbackslash mathbb\{Z\}\$.

\subsubsection{4.1 Content Addressing}\label{content-addressing}

Every atom, morpheme, word, and verse receives a SHA-256 hash as its
address. Addressing is compositional and bottom-up:

\$\$\textbackslash text\{atoms\}
\textbackslash;\textbackslash longrightarrow\textbackslash;
\textbackslash text\{morphemes\}
\textbackslash;\textbackslash longrightarrow\textbackslash;
\textbackslash text\{words\}
\textbackslash;\textbackslash longrightarrow\textbackslash;
\textbackslash text\{verses\}\$\$

A word\textquotesingle s address is determined entirely by its
constituent atoms, making identity intrinsic to content.

\subsubsection{4.2 Dihedral Group
Structure}\label{dihedral-group-structure}

The ring carries a dihedral group \$D\_\{2\^{}\{256\}\}\$ generated by
two involutions:

\$\$\textbackslash operatorname\{bnot\}(a) = 2\^{}\{256\} - 1 - a
\textbackslash qquad\textbackslash text\{(bitwise NOT, reflection)\}\$\$

\$\$\textbackslash operatorname\{neg\}(a) = 2\^{}\{256\} - a
\textbackslash qquad\textbackslash text\{(arithmetic negation,
rotation)\}\$\$

Every address belongs to a 4-element orbit:
\$\textbackslash\{a,\textbackslash;\textbackslash operatorname\{bnot\}(a),\textbackslash;\textbackslash operatorname\{neg\}(a),\textbackslash;\textbackslash operatorname\{bnot\}(\textbackslash operatorname\{neg\}(a))\textbackslash\}\$.

\subsubsection{4.3 Metrics}\label{metrics}

Metric

Definition

Measures

Ring distance

\$\textbackslash min(\textbar a-b\textbar,\textbackslash; 2\^{}\{256\} -
\textbar a-b\textbar)\$

Algebraic proximity

Hamming distance

Bit differences in hex

Structural similarity

Incompatibility

\$\textbackslash operatorname\{hamming\}(a,\textbackslash operatorname\{bnot\}(b))\$

Complementarity

Table 4. UOR distance metrics.

\subsubsection{4.4 Fiber Decomposition}\label{fiber-decomposition}

The \textbf{2-adic valuation} \$v\_2(a)\$ (number of trailing zero bits)
assigns each address to a stratum. This creates a natural fibration:
stratum 0 (odd addresses) is the generic fiber, and higher strata are
increasingly structured.

\subsection{5. Integration}\label{integration}

The three frameworks serve different analytical purposes:

\begin{enumerate}
\tightlist
\item
  \textbf{Root Space} answers: \emph{What does this root do across the
  Quran? Is this usage normal or exceptional?}
\item
  \textbf{su(2) Gauge Field} answers: \emph{Is this
  root\textquotesingle s meaning conserved or evolving? Where are the
  tension points?}
\item
  \textbf{UOR} answers: \emph{What is the intrinsic identity of this
  textual unit? How do units relate compositionally?}
\end{enumerate}

A typical analysis flow:

\begin{enumerate}
\tightlist
\item
  Build root vectors for a passage {(root space)}
\item
  Compute per-word dynamics: weight, anomaly, drift, coherence {(verse
  dynamics)}
\item
  Detect boundaries and emphasis patterns {(verse dynamics)}
\item
  For roots of interest, compute holonomic resonance {(gauge field)}
\item
  Identify high-deviation instances as falsification targets
  {(Lyapunov)}
\item
  Track compositional identity through the text hierarchy {(UOR)}
\end{enumerate}

\subsubsection{5.1 Tool Mapping}\label{tool-mapping}

Tool

Math used

\texttt{analyze\_verse\_emphasis}

Root space + verse dynamics

\texttt{detect\_boundaries}

Verse coherence (phase-lock)

\texttt{measure\_phase\_lock}

Verse coherence between two verses

\texttt{compute\_passage\_drift}

Per-word drift energy

\texttt{compute\_root\_resonance}

\$\textbackslash mathfrak\{su\}(2)\$ holonomy

\texttt{verify\_root\_concordance}

Lyapunov deviation

\texttt{root\_distance}

Jensen--Shannon divergence

\texttt{compare\_with\_traditional}

Verse dynamics emphasis profile

Table 5. Mapping of MCP tools to mathematical frameworks.

\end{document}
